\chapter{A Benchtop Sandbox: Simple Systems for Practising Experimental Design}
\label{ch:benchtop}

\begin{quote}\itshape
``For the things we have to learn before we can do them, we learn by doing
them.'' --- Aristotle, \emph{Nicomachean Ethics}.
\end{quote}

\noindent
The models of the preceding chapters are faithful, and for that reason they are
complicated. A chromatographic band is a coupled system of partial differential
equations (Chapter~\ref{ch:models}); a fermentation is a stochastic,
indirectly-observed community of living cells (Chapter~\ref{ch:ferment}). When a
design method succeeds or fails on such a system, it can be hard to say whether
the lesson lies in the \emph{method} or in the \emph{mechanism}. This chapter
steps deliberately down in complexity. It assembles four \emph{benchtop}
systems---each an experiment a student could run in an afternoon, each governed
by one or two equations you can hold in your head---and uses them as a sandbox in
which the behaviour of an experimental design is laid bare.

The four are chosen so that between them they exhibit the response \emph{shapes}
that most often decide which design to reach for:

\begin{itemize}
\item \textbf{Pressure drop in a pipe} (Section~\ref{sec:bt-pipe}) is a pure
  \emph{power law}. It is the ideal system on which to see why a logarithmic
  transform turns a tangle of multiplicative effects into clean, additive ones
  that a two-level \emph{factorial} estimates exactly.
\item The \textbf{falling-ball viscometer} (Section~\ref{sec:bt-ball}) is a
  \emph{calibration} problem: the quantity we want (viscosity) is not the
  quantity we measure (a fall time), and the design question is \emph{where} to
  operate so the inversion is trustworthy---a lesson in \emph{parameter
  estimation} and measurement bias.
\item \textbf{Back-extrusion of a set yogurt} (Section~\ref{sec:bt-yogurt}) has a
  \emph{threshold}: the yield stress. Its response surface is curved and its
  ``pass/fail'' texture specification is naturally probabilistic, exercising
  \emph{response-surface} and \emph{generalized-linear} models---and tying back
  to the dairy fermentation of Chapter~\ref{ch:ferment}.
\item The \textbf{hydrolysis of methyl acetate} (Section~\ref{sec:bt-ester}) is a
  slow \emph{reversible reaction}. Because it takes hours to approach
  equilibrium, \emph{time} joins temperature and catalyst as a design factor, and
  the system cleanly separates an effect on \emph{rate} (the catalyst) from an
  effect on \emph{position of equilibrium} (the temperature).
\end{itemize}

Every system lives in \code{vlab\_doe.models.benchtop}, and every figure in
this chapter is generated by running that code. Section~\ref{sec:bt-choosing}
closes with a decision table: read the shape of your response, and let it name
the method.

%==============================================================================
\section{Pressure drop in a pipe: power laws and the log transform}
\label{sec:bt-pipe}

A pump pushes water through a straight horizontal tube of length $L$ and internal
diameter $d$, and we read the pressure drop $\Delta P$ across it on a manometer at
a series of volumetric flow rates $Q$. Nothing could be simpler to perform; yet
the response carries almost every feature a first course in DoE wants to teach.

\subsection{The mechanism}

With mean velocity $v = Q/A$ over the cross-section $A = \pi d^2/4$, and Reynolds
number $\mathrm{Re} = \rho v d/\mu$, the Darcy--Weisbach equation writes the
pressure drop through a dimensionless friction factor $f$,
%
\begin{equation}
\Delta P = f\,\frac{L}{d}\,\frac{\rho v^2}{2}.
\label{eq:darcy}
\end{equation}
%
In \emph{laminar} flow ($\mathrm{Re}\lesssim 2300$) the friction factor is
exactly $f = 64/\mathrm{Re}$, and \eqref{eq:darcy} collapses to the
\emph{Hagen--Poiseuille} law,
%
\begin{equation}
\Delta P = \frac{128\,\mu L Q}{\pi d^4},
\label{eq:hagen}
\end{equation}
%
a response \emph{linear} in $Q$ and $L$ but falling as the \emph{fourth power} of
the diameter. In \emph{turbulent} flow ($\mathrm{Re}\gtrsim 4000$) a smooth-pipe
Blasius correlation $f = 0.316\,\mathrm{Re}^{-1/4}$ gives $\Delta P \propto
Q^{1.75}$. The package (\code{pipe\_flow.py}) evaluates \eqref{eq:darcy} with a
friction factor that interpolates smoothly between the two branches, so the
characteristic is continuous through the transition (Figure~\ref{fig:bt-pipe}).
The water properties $\mu(T)$ and $\rho(T)$ are the Vogel and Kell correlations,
so \emph{temperature} enters as a genuine---if gentle---fifth factor.

\begin{figure}[t]
\centering
\includegraphics[width=\linewidth]{benchtop_pipe.png}
\caption[The pipe pressure-drop characteristic]{The pressure-drop characteristic
of a $1$~m, $4$~mm pipe carrying water at $20^\circ$C, from
\code{pipe\_flow.flow\_curve}. \emph{Left:} $\Delta P$ rises linearly with flow
in the laminar regime---exactly Hagen--Poiseuille, Eq.~\eqref{eq:hagen}---then
steepens to the $\sim\!Q^{1.75}$ turbulent law past the transition. \emph{Right:}
the friction factor is $64/\mathrm{Re}$ in the laminar branch and joins the
Blasius line above the (shaded) transition band.}
\label{fig:bt-pipe}
\end{figure}

\begin{databox}{the pipe experiment}
The apparatus is a peristaltic or gear pump, a length of precision-bore tubing,
and a differential pressure transducer. Each ``run'' sets a flow rate and reads
the steady pressure drop; a full characteristic is a dozen flow rates in a few
minutes. Replication is cheap, so the \emph{pure-error} estimate that anchors an
ANOVA is easy to obtain. The controllable factors are the flow rate $Q$ (pump
speed), the tube geometry $L$ and $d$ (swap the tubing), and the temperature $T$
(a water bath). The response $\Delta P$ spans several orders of magnitude across a
realistic design region---which is precisely why it must be modelled on a log
scale.
\end{databox}

\subsection{The design lesson: a factorial in the logarithms}

Because \eqref{eq:hagen} is a \emph{product} of powers, taking logarithms makes
it \emph{additive}:
%
\begin{equation}
\log \Delta P = \text{const} + 1\cdot\log Q + 1\cdot\log L - 4\cdot\log d + \dots
\label{eq:loglinear}
\end{equation}
%
On the log scale the mechanistic exponents \emph{are} the main-effect
coefficients, and the interactions that clutter the raw-scale response vanish. A
two-level full factorial in $(\log Q,\log L,\log d)$---eight runs plus a few
centre points---therefore recovers the physics exactly. Running the package's
\code{full\_factorial} through \code{pipe\_flow.pressure\_drop} and fitting with
\code{fit\_response\_model} returns slopes of $+1.00$, $+1.00$, and $-4.00$ to
two decimal places, with $R^2$ indistinguishable from one
(Figure~\ref{fig:bt-pipe-effects}). The same eight runs fitted on the \emph{raw}
scale would demand large two- and three-factor interaction terms and still fit
poorly---the classic symptom of a response that lives on the wrong scale.

\begin{figure}[t]
\centering
\includegraphics[width=\linewidth]{benchtop_pipe_effects.png}
\caption[A factorial recovers the power-law exponents]{A $2^3$ factorial on the
log-transformed factors, evaluated through the pipe model. \emph{Left:} the
fitted log--log slopes coincide with the mechanistic exponents $(+1,+1,-4)$ of
Hagen--Poiseuille (black diamonds). \emph{Right:} predicted versus observed
$\log\Delta P$ falls on the diagonal ($R^2\approx 1$). The lesson is not about
pipes: it is that the \emph{scale} on which a factorial is analysed can turn an
intractable interaction structure into three clean numbers.}
\label{fig:bt-pipe-effects}
\end{figure}

\begin{keybox}{Reading the scale off the mechanism}
Whenever a response is a product of powers of the factors---and rate laws,
transport correlations, and cost models very often are---its natural home is the
log scale, where multiplicative effects become additive main effects and a small
factorial suffices. When a factorial ``needs'' large interactions to fit, suspect
the scale before you enlarge the design.
\end{keybox}

%==============================================================================
\section{The falling-ball viscometer: calibration and the cost of the wrong operating point}
\label{sec:bt-ball}

Our second system inverts the logic of the first. There, we knew the mechanism
and used the design to \emph{confirm} it; here we know the mechanism but use it to
\emph{measure} an unknown---the viscosity of a fluid---and the design question is
where to operate so the measurement is unbiased.

\subsection{The mechanism}

A dense sphere of diameter $d$ and density $\rho_s$ is dropped into a vertical
tube of internal diameter $D$ filled with the test fluid ($\rho_f$, $\mu$), and we
time its fall over a marked distance. At terminal velocity the buoyant weight
balances the drag; in the creeping-flow (Stokes) limit this solves in closed
form,
%
\begin{equation}
v_{\mathrm{St}} = \frac{g\,d^2(\rho_s-\rho_f)}{18\,\mu},
\qquad
\mathrm{Re} = \frac{\rho_f v d}{\mu},
\label{eq:stokes}
\end{equation}
%
so a single fall time $t = L/v$ yields $\mu$. Two corrections turn the textbook
formula into a realistic instrument, and both are built into
\code{falling\_ball.py}. A \emph{wall} slows the ball: the Ladenburg--Faxén factor
$K_w = 1 - 2.104\,\beta + 2.09\,\beta^3 - 0.95\,\beta^5$, with $\beta = d/D$,
multiplies $v_{\mathrm{St}}$. And \emph{inertia} adds drag beyond Stokes': the
package solves the force balance with the Schiller--Naumann coefficient
$C_D = (24/\mathrm{Re})(1 + 0.15\,\mathrm{Re}^{0.687})$, so the true terminal
velocity falls below \eqref{eq:stokes} as $\mathrm{Re}$ grows.

\subsection{The design lesson: choose the operating point, then estimate}

An experimenter who reads the fall time and inverts the \emph{Stokes} formula
\eqref{eq:stokes}---as the classic instrument does---commits a bias whenever the
ball is not in creeping flow. Because $\mathrm{Re}$ grows with the ball's size and
excess density, that bias is a function of the \emph{operating point}
(Figure~\ref{fig:bt-ball}, right). The design task is therefore twofold: pick a
ball and tube that keep $\mathrm{Re}\ll 1$ (so the cheap Stokes inversion is
valid) \emph{and} keep the fall time long enough to time precisely but short
enough to be practical. This is a genuine multi-objective design over
$(d,\rho_s,D)$, and the package lets one run it as a factorial through
\code{fall\_time} and score each corner by both its fall time and its
Stokes-inversion error.

\begin{figure}[t]
\centering
\includegraphics[width=\linewidth]{benchtop_falling_ball.png}
\caption[Falling-ball calibration and its bias]{\emph{Left:} the calibration
curve---fall time over $0.1$~m versus true viscosity for a $2$~mm steel ball,
from \code{fall\_time}---is a clean straight line on log--log axes over two
decades of viscosity, the working range of the instrument. \emph{Right:} the
error incurred by inverting with the naive Stokes formula, as a function of the
particle Reynolds number: negligible for a small, slow ball, but growing to
several percent as a larger ball leaves creeping flow. The design choice of ball
and tube \emph{is} the choice of how trustworthy the reading will be.}
\label{fig:bt-ball}
\end{figure}

\code{infer\_viscosity} exposes both inversions: the analytic wall-corrected
Stokes reading (\code{stokes\_only=True}), which carries the bias, and a numerical
inversion of the full drag law (\code{stokes\_only=False}), which recovers the
viscosity used to generate the data to four significant figures. Running a
factorial of ball choices through the instrument and inverting each with the full
model returns the same $\mu$ from every corner---a reassuring \emph{internal
consistency} check that a good calibration design should pass.

\begin{limitbox}{the falling-ball model}
The model assumes a \emph{single, isolated} sphere at its terminal velocity in a
Newtonian fluid: it ignores the entry length over which the ball accelerates
(short for dense balls in viscous fluids), end effects near the tube's top and
bottom, and any non-Newtonian character of the fluid. The wall correlation is
accurate only for slender balls ($\beta \lesssim 0.6$), and the drag law is
fitted for $\mathrm{Re}\lesssim 10^3$; the package clips $\beta$ and brackets the
root accordingly. For a genuinely non-Newtonian fluid, the back-extrusion model
of the next section is the appropriate tool.
\end{limitbox}

%==============================================================================
\section{Back-extrusion of a set yogurt: a threshold in the response}
\label{sec:bt-yogurt}

The third system measures the \emph{texture} of a semi-solid---set yogurt,
custard, labneh---and connects this chapter to the dairy fermentation of
Chapter~\ref{ch:ferment}: the very product whose making we modelled there, we now
put under a probe.

\subsection{The mechanism}

A texture analyser pushes a flat cylindrical probe of radius $R_p$ down at
constant speed $V$ into a cup of radius $R_c$ filled to depth $H$, and records the
force. The displaced product flows up through the thin annular gap. Yogurt is
\emph{shear-thinning with a yield stress}, described by the Herschel--Bulkley law
%
\begin{equation}
\tau = \tau_0 + K\,\dot\gamma^{\,n},
\label{eq:hb}
\end{equation}
%
with consistency $K$, flow index $n<1$, and yield stress $\tau_0$---the three
numbers a fermentation actually changes. Modelling the annulus as a thin slit,
\code{back\_extrusion.py} sums a yield contribution and a power-law (consistency)
contribution to the pressure gradient the probe must generate, and multiplies by
the probe face to obtain the force. Warmer product is softer: $K$ and $\tau_0$
carry an Arrhenius temperature factor, so a gel measured at $5^\circ$C reads
firmer than the same gel at room temperature (Figure~\ref{fig:bt-yogurt}, right).

\begin{figure}[t]
\centering
\includegraphics[width=\linewidth]{benchtop_back_extrusion.png}
\caption[Back-extrusion texture curves]{Texture responses from
\code{back\_extrusion.peak\_force}. \emph{Left:} peak force versus probe speed for
three yield stresses. Because $n=0.4<1$ the curves are \emph{concave}---force
grows far slower than linearly with speed---and the yield stress sets a
speed-independent offset that separates a ``drinkable'' from a ``firm-set''
product. \emph{Right:} the same gel softens markedly as it warms, a reminder that
\emph{sample temperature} is a factor to be controlled, not assumed.}
\label{fig:bt-yogurt}
\end{figure}

\subsection{The design lesson: curvature and a probabilistic specification}

Two features make this the natural home for the response-surface and
generalized-linear methods of Chapter~\ref{ch:design}. First, the response is
genuinely \emph{curved} in the operating factors---speed enters through a
power $n$, temperature through an exponential---so a two-level factorial is not
enough: a three-level design (or a central-composite) is needed to estimate the
quadratic terms, and a space-filling \emph{Latin hypercube}
(\code{latin\_hypercube}) is a good way to explore the $(V,H,R_p,T)$ box before
committing to a surface. Second, the industrial question is rarely ``what is the
force?'' but ``does the batch fall within the texture \emph{specification}?''---a
pass/fail outcome. Thresholding the peak force at a target window turns the design
into a \emph{probabilistic design space} problem, fitted with
\code{fit\_glm\_response} exactly as the purity and breakthrough specifications
were in Chapter~\ref{ch:design}. The yield stress, meanwhile, gives the surface a
\emph{floor}: below a critical speed the force is dominated by $\tau_0$ and barely
responds to $V$ at all---a plateau that a purely linear model would badly
misread.

\begin{appbox}{quality control of cultured dairy}
Back-extrusion and its cousin, the penetration test, are the workhorse
instrumental measurements of dairy texture. A yogurt line specifies a firmness
window; drift outside it signals an over- or under-set batch, a fault traceable
to the fermentation (Chapter~\ref{ch:ferment})---the wrong strain blend,
temperature, or set-point pH. Coupling the fermentation model's output rheology
to this measurement closes the loop from \emph{process} parameters to a
\emph{product} attribute a consumer actually feels, which is the essence of a
quality-by-design specification.
\end{appbox}

%==============================================================================
\section{Hydrolysis of methyl acetate: rate, equilibrium, and time as a factor}
\label{sec:bt-ester}

The final system is a chemical reaction slow enough that \emph{when} you look is
part of the design.

\subsection{The mechanism}

Methyl acetate reacts with water, under acid catalysis, to acetic acid and
methanol:
%
\begin{equation}
\mathrm{CH_3COOCH_3} + \mathrm{H_2O}
\;\underset{k_r}{\overset{k_f}{\rightleftharpoons}}\;
\mathrm{CH_3COOH} + \mathrm{CH_3OH}.
\label{eq:ester}
\end{equation}
%
The reaction is \emph{reversible}, reaching an equilibrium short of full
conversion, and slow---hours at room temperature. The hydrogen ion is a
\emph{catalyst}: it multiplies the forward and reverse rates equally, so with the
reaction extent $x$ (mol/L),
%
\begin{equation}
\dv{x}{t} = c_{\mathrm{H^+}}\big[\,k_f(E_0-x)(W_0-x) - k_r(A_0+x)(M_0+x)\,\big],
\label{eq:esterode}
\end{equation}
%
where $k_f, k_r$ are Arrhenius in temperature. Two consequences organise the
design, and \code{ester\_hydrolysis.py} is tuned to display both. Because the
catalyst multiplies \emph{both} directions, its ratio---the equilibrium constant
$K_{\mathrm{eq}} = k_f/k_r$---is independent of $c_{\mathrm{H^+}}$: \textbf{acid
speeds the approach to equilibrium but does not move it}
(Figure~\ref{fig:bt-ester}, left). Temperature, by contrast, changes $k_f$ and
$k_r$ \emph{unequally}, so it moves the rate \emph{and}, through van 't Hoff, the
equilibrium itself; with the package's mildly endothermic defaults, warmer runs
convert more ester (Figure~\ref{fig:bt-ester}, right). An operating constraint
completes the picture: methyl acetate boils at $56.9^\circ$C, so the temperature
factor has a hard ceiling, and \code{EsterHydrolysisConfig} refuses to be built
above it.

\begin{figure}[t]
\centering
\includegraphics[width=\linewidth]{benchtop_ester.png}
\caption[Ester-hydrolysis kinetics and equilibrium]{\emph{Left:} conversion
trajectories from \code{ester\_hydrolysis.simulate} at three acid
concentrations. More catalyst reaches the plateau sooner, but \emph{all three
share the same equilibrium} (dotted line): the catalyst is a rate lever, not an
equilibrium lever. \emph{Right:} the equilibrium conversion rises with
temperature (the reaction is endothermic as parameterised), up to the hard
$56.9^\circ$C boiling-point ceiling of methyl acetate.}
\label{fig:bt-ester}
\end{figure}

\subsection{The design lesson: separating levers, with time in the design}

Here the design factors are temperature $T$, catalyst $c_{\mathrm{H^+}}$ (or pH,
via \code{from\_ph}), the water-to-ester ratio, and---uniquely in this
chapter---the \emph{sampling time} $t$. A factorial in $(T, c_{\mathrm{H^+}}, t)$
with conversion as the response, evaluated through \code{conversion\_at},
produces an instructive ANOVA: a strong $c_{\mathrm{H^+}}\times t$
\emph{interaction} (more acid matters more early, and not at all once equilibrium
is reached), a main effect of $T$ that acts through \emph{both} rate and
equilibrium, and---if the sampling times straddle the approach to
equilibrium---pronounced curvature in $t$. Plotting conversion over the
$(T,t)$ plane (Figure~\ref{fig:bt-ester-surface}) makes the two regimes visible
at once: a rate-limited rise at short times, flattening onto the
temperature-dependent equilibrium ridge. It is a compact illustration of why a
kinetic response demands that \emph{time} be treated as a designed factor rather
than a nuisance held fixed.

\begin{figure}[t]
\centering
\includegraphics[width=0.72\linewidth]{benchtop_ester_surface.png}
\caption[The conversion response surface]{Ester conversion over temperature and
reaction time at fixed acid ($c_{\mathrm{H^+}}=0.1$~M), from
\code{conversion\_at}. Short times are rate-limited (steep in $t$); long times
flatten onto the equilibrium ridge, which itself creeps upward with temperature.
A design that sampled at a single time would see only a slice of this surface---
and could confound a rate effect with an equilibrium one.}
\label{fig:bt-ester-surface}
\end{figure}

\begin{keybox}{One factor, two mechanisms}
The ester system warns against a common misreading: a single factor may act
through more than one mechanism. Temperature here moves \emph{both} the rate and
the equilibrium, and only a design that samples across time can tell the two
apart. When an effect seems to change sign or magnitude with the observation
window, suspect that the factor is pulling more than one lever.
\end{keybox}

%==============================================================================
\section{A worked example: a student's design--run--analyse cycle}
\label{sec:bt-worked}

To see the pieces of this book move together, follow a student---call her
Maya---given a single brief on the ester system of the previous section:
\emph{find operating conditions that hold the hydrolysis conversion inside a
$0.18$--$0.30$ window, and report how far each factor may drift while staying in
spec.} She has no wet lab today; her ``reactor'' is the benchtop model, and her
whole workflow is a few dozen lines against the package. It is worth watching not
because the chemistry is deep---it is not---but because the \emph{procedure} is
exactly the one she will later run on a chromatography skid, where each
experiment costs a day and a column.

\subsection{Framing the design}

Maya lists three controllable factors and honest ranges: temperature
$25$--$50^\circ$C (below the $56.9^\circ$C boiling ceiling that
\code{EsterHydrolysisConfig} enforces), catalyst $c_{\mathrm{H^+}}=0.05$--$0.40$
mol/L, and sampling time $0.5$--$4$~h. Her response is the fractional conversion.
Crucially, she does not treat the model as an oracle of perfect numbers: she wraps
its output in the \code{perturbation} module, adding the $2\%$ proportional and
$0.004$ additive measurement noise of a real assay, so every ``run'' is a
\emph{virtual experiment} with the scatter a titration would actually show. A
fixed seed (\code{make\_rng}) keeps the study reproducible.

\subsection{Stage one---screen, and read the diagnostics}

She starts cheap. \code{full\_factorial} builds a two-level $2^3$ design and, on
her insistence, \emph{four replicated centre points}---the small insurance that
will shortly earn its keep. \code{run\_design} pushes all twelve conditions
through the noisy model, and \code{fit\_response\_model} fits a first-order model
with two-factor interactions and returns its Type-II ANOVA.

The screen is, on its face, a disappointment. The model explains only
$R^2\approx0.74$ of the variation, and \emph{not one} term clears $p<0.05$:
catalyst is the largest effect but its $F$-test is marginal, and the interactions
are buried in noise. A na\"ive reading would be ``nothing is significant; collect
more replicates.'' That reading is wrong, and the centre points say why. Their
mean conversion is $0.32$, while the mean over the factorial corners is only
$0.22$: the response \emph{bulges} above the flat plane a two-level design can
fit, by fully $0.10$---a curvature signal several times the noise. The poor $R^2$
is not a lack of \emph{signal} but a lack of \emph{model}; the first-order surface
is simply the wrong shape (Figure~\ref{fig:bt-workflow}, left, where the
replicated centre points sit far off the diagonal because the model can only
predict their value as the average of the corners). The screen has done its real
job: it has told her to change the model, not to gather more data.

\subsection{Stage two---escalate to a response surface}

Curvature demands quadratic terms, so Maya augments to a three-level design
(\code{full\_factorial(\dots, levels=3)}) and refits with
\code{quadratic=True}. The picture snaps into focus: $R^2\approx0.93$, and the
ANOVA now speaks clearly. Time and catalyst are the dominant main effects
($p<10^{-3}$); both carry significant \emph{quadratic} terms---the plateau as
conversion approaches equilibrium; and the $c_{\mathrm{H^+}}\times t$ interaction
is real ($p\approx0.007$), the mathematical echo of the lesson from
Section~\ref{sec:bt-ester} that more acid matters early and not at all once
equilibrium is reached. Temperature, tellingly, is weak as a \emph{main} effect
yet significant through its interactions with both catalyst and time---precisely
the ``one factor, two mechanisms'' signature flagged above, now falling out of the
ANOVA table on its own.

\begin{figure}[t]
\centering
\includegraphics[width=\linewidth]{benchtop_student_workflow.png}
\caption[A student's design--run--analyse cycle]{The full workflow on the ester
system, produced entirely by the package. \emph{Left:} predicted versus observed
conversion. The two-level screen (blue) fits the corners passably but its four
replicated centre points (dark squares) fall well off the diagonal---it predicts
their value as the corner average and misses the real, higher conversion, the
unmistakable signature of curvature. The three-level response-surface model
(green) hugs the diagonal ($R^2\approx0.93$). \emph{Right:} the fitted conversion
surface over the two dominant factors at $37.5^\circ$C, with the design points
overlaid and the $0.18$--$0.30$ specification window hatched---the proven
acceptable region Maya was asked to deliver.}
\label{fig:bt-workflow}
\end{figure}

\subsection{Stage three---turn the surface into an answer}

A fitted surface is not yet an answer; the brief asked for \emph{ranges}. Maya
calls \code{proven\_acceptable\_ranges} with the specification
$(0.18,0.30)$. Holding each factor in turn at the centre of the others, it inverts
the fitted model to report how far each may move before the prediction leaves the
window. The binding constraint is the upper bound---conversion presses against its
temperature-set equilibrium ceiling---so the acceptable region is a modest-acid,
short-time corner: catalyst up to $\approx0.20$~mol/L and time up to
$\approx2.0$~h keep her inside spec, with temperature admissible up to
$\approx35^\circ$C at the centre point. That corner is exactly the hatched region
of Figure~\ref{fig:bt-workflow} (right), and it is the deliverable: not a single
``optimal'' recipe but a \emph{window}, the language in which robust processes are
actually specified.

\begin{databox}{the virtual experiment}
Every number in this section came from the same three ingredients the rest of the
book uses: a mechanistic model (\code{ester\_hydrolysis}), a noise layer
(\code{perturbation.add\_measurement\_noise}) that turns clean truth into a
realistic measurement, and the design/analysis code
(\code{full\_factorial}, \code{run\_design}, \code{fit\_response\_model},
\code{proven\_acceptable\_ranges}). Because the ``truth'' is known here, the study
is also a \emph{check on the methods}: one can ask whether the recovered surface
and window match what the model would give without noise---a luxury no wet lab
ever has, and the reason a virtual sandbox is worth building.
\end{databox}

\subsection{What the same data could do next}

Maya's cycle used the classical factorial and response-surface tools, but the very
same runs feed the rest of the toolkit developed in Chapter~\ref{ch:design}. Had
her specification been a pass/fail texture or purity outcome rather than a
continuous conversion, \code{fit\_glm\_response} would have fitted a logistic
\emph{probabilistic} design space over the same factors. Had the region been
higher-dimensional or expensive to sample, \code{latin\_hypercube} would have
space-filled it more efficiently than a grid, and the Bayesian-optimisation loop of
Section~\ref{sec:bo} would have driven successive runs toward an optimum instead of
mapping the whole surface. And had the goal been to \emph{estimate the kinetics}
rather than to operate the reaction, the package's inverse-modelling and
uncertainty tools (the \code{uq} module) would fit $k_f$, $k_r$, and their
activation energies to the same conversion data, with credible intervals. The
sandbox does not replace those methods; it is the place to build the judgement
that deploys them.

%==============================================================================
\section{Choosing a method by the shape of the response}
\label{sec:bt-choosing}

The four systems were chosen to make a single point: the right experimental
design is dictated less by the field of application than by the \emph{shape} of
the response. Table~\ref{tab:bt-summary} distils the chapter into that mapping.

\begin{table}[t]
\centering
\caption[Benchtop systems and the methods they teach]{The four benchtop systems,
the response feature each foregrounds, and the design method it motivates.}
\label{tab:bt-summary}
\small
\begin{tabular}{@{}p{0.24\linewidth}p{0.30\linewidth}p{0.38\linewidth}@{}}
\toprule
\textbf{System} & \textbf{Response shape} & \textbf{Method it teaches} \\
\midrule
Pipe pressure drop & Product of powers; regime change &
Two-level \emph{factorial} on \emph{log}-transformed factors; effect estimation
and the choice of analysis scale. \\
Falling-ball viscometer & Monotone, but the target is inferred, not measured &
\emph{Parameter estimation} and calibration; designing the operating point to
control bias. \\
Yogurt back-extrusion & Curved, with a yield-stress threshold and a pass/fail spec &
\emph{Response-surface} (3-level / LHS) and \emph{generalized-linear} models for a
probabilistic design space. \\
Methyl-acetate hydrolysis & Kinetic approach to a shifting equilibrium &
Designs with \emph{time} as a factor; separating rate from equilibrium via
interactions and curvature. \\
\bottomrule
\end{tabular}
\end{table}

None of these systems needs a supercomputer, a chromatography skid, or a
bioreactor; each can be simulated in milliseconds and, in a teaching laboratory,
performed in an afternoon. That is the point. Stripped of mechanistic
complication, the designs of Chapter~\ref{ch:design} show their working, and the
intuition earned here---read the response's shape, choose the scale, put the
right factors (including time) in the design, and mind the difference between what
you measure and what you want---transfers directly back to the harder problems
with which the monograph began.

\begin{keybox}{the benchtop sandbox}
Four simple systems---a pipe, a falling ball, a probe in a gel, a slow
reaction---span the response shapes that decide a design: power law, calibration,
threshold, and kinetics. Practising the methods where the mechanism is
transparent builds the judgement needed where it is not.
\end{keybox}
